We separately benchmarked Qwen2.5-0.5B on the Intel Core i7-10700 and AMD Ryzen 7 8700G using the same Torch and Transformers versions, prompt, and inference procedure, without cache probes.
After two warmup requests, we measured ten requests with no prefix-cache reuse.
The hybrid configurations offload all 168 transformer linear projections across 24 decoder blocks to synchronous, single-AIE-core NPU kernels; attention, nonlinear operations, and the output head remain on the CPU.
DQ uses an activation threshold of 6.0 and CPU-side outlier gathering and correction.
Table~\ref{tab:architecture-performance} reports implementation-specific performance: DQ lowers the measured memory footprint in all three execution settings, but reduces TTFT only for the AMD CPU in these tests.
The Intel INT8 backend falls back to FP32 matrix multiplication, whereas the AMD CPU uses a native integer kernel.
The hybrid implementation includes transfer and dispatch overheads and does not outperform the CPU baselines; these measurements do not represent optimized full-array NPU performance.
